\section*{Original bearing geometric statements (pp.9:6,9:12,9:14--15)}
Given $A,B\subseteq\mathbb R^n$, their Minkowski sum is defined as
\[A+B:=\{a+b:a\in A,b\in B\}.\]
If $P$ and $Q$ are polytopes then $P+Q$ is also a polytope. In two-dimensions, Minkowski sum has nice properties. Let $P$ be a polytope in $\mathbb R^2$ with vertices $v_1,\ldots,v_k$ where $k>1$. We can assume they are ordered so that $P$ lies in the left closed half plane determined by the line going from $v_i$ to $v_{i+1}$ for $i<k$, and similarly for $v_k$ and $v_1$. Let $E(P)$ be the collection of unit vectors in the direction of these $k$ edges. That is, vectors of the form $(v_{i+1}-v_i)/\|v_{i+1}-v_i\|$ for $i<k$, and $(v_1-v_k)/\|v_1-v_k\|$. If $|\vertset(P)|\le1$ then $E(P):=\emptyset$.
\begin{sourceresult}{Lemma 8}
Let $P_1,\ldots,P_r$ be non-empty polytopes in $\mathbb R^2$. Then $E(P_1+\cdots+P_r)=\bigcup_{i=1}^r E(P_i)$. Consequently, $|\vertset(P_1+\cdots+P_r)|\le\sum_{i=1}^r|\vertset(P_i)|$. The latter holds for empty $P_i$'s as well.
\end{sourceresult}
The lemma is folklore. It can be inferred from Chapter 13.3 in [11], and we give only an outline of proof.
\begin{proof}[Proof sketch of Lemma 8]
Given a non-empty polytope $P$ and $w\in\mathbb R^2$, let $P^w:=\{x\in P:\langle x,w\rangle=\max_{z\in P}\langle z,w\rangle\}$ be the set of extreme of points of $P$ in the direction $w$. It is either a vertex or an edge of $P$. For a pair of polytopes we have $(P_1+P_2)^w=P_1^w+P_2^w$. Every edge of $P_1$ yields an edge of $P_1+P_2$ with the same direction. Conversely, every edge of $P_1+P_2$ comes from one of $P_1$ or $P_2$.
\end{proof}

\begin{sourceresult}{Lemma 28}
Let $f$ be a polynomial over an infinite field $\mathbb F$. Assume that $\sigma(\Newt(f))=k$. Then there exists a bivariate Laurent polynomial $g\in\mathbb F(y_1,y_2,y_1^{-1},y_2^{-1})$ which is an h.p.-projection of $f$ so that $\Newt(g)$ has $k$ vertices. Moreover, if $f$ is a homogeneous polynomial then $g$ is a polynomial. If $\mathbb F=\mathbb R$, then the coefficients in the projection can be assumed positive.
\end{sourceresult}
\begin{proof}
Let $L(z)=Az$ with $A\in\mathbb R^{2\times n}$ be a linear map so that
\[|\vertset(L(\Newt(f)))|=k.\]
We can assume that the entries of $A$ are rational, because a small perturbation of $A$ cannot decrease $|\vertset(L(\Newt(f)))|$. Now, we can assume that the entries of $A$ are integers, because we can multiply $A$ by a suitable integer.
Moreover, when $f$ is homogeneous of degree $d$, increasing all entries of $A$ by $\lambda$ corresponds to shifting $L(\Newt(f))$ by $(\lambda d,\lambda d)$, which does not change the number of vertices. Hence, in this case, $A$ can be taken with non-negative integer entries.
Let us now consider a h.p.-projection $\pi$ with $\pi(x_i)=a_i y^{L(e_i)}$. It follows that $\supp(\pi(f))\subseteq L(\supp(f))$. Now, we claim that we can choose the coefficients $a_i$ so that equality holds. This can be seen as follows. Given $\alpha\in\supp(f)$, the coefficient of $y^{L(\alpha)}$ in $\pi(f)$ is a non-zero polynomial $h_\alpha$ in the coefficients of $\pi$. Hence, if $\mathbb F$ is infinite, there exist non-zero coefficients for which $h_\alpha$ is non zero for every $\alpha\in\supp(f)$. If $\mathbb F=\mathbb R$, they can be taken positive.
\end{proof}

\begin{sourceresult}{Lemma 37}
Let $A,B\subseteq\mathbb R^2$ be non-empty sets such that $A+B$ is a convexly independent set with $|A|\ge|B|$. Then either $|A|\le2$ or $|B|\le1$.
\end{sourceresult}
\begin{proof}
For the sake of contradiction, assume that $A+B$ is convexly independent, $|A|\ge3$ and $|B|\ge2$. By Lemma 8, the convex hull of $A+B$ has at most $|A|+|B|$ vertices. By the size assumption, there exist $a_1\ne a_2\in A$ and $b_1\ne b_2\in B$ with $a_1+b_1=a_2+b_2$. The point $a_1+b_1$ is the average of $a_1+b_2$ and $a_2+b_1$ and it is distinct from them, a contradiction.
\end{proof}
